\documentclass[12pt,a4paper]{article}
\usepackage[utf8]{inputenc}
\usepackage[T1]{fontenc}
\usepackage{amsmath}
\usepackage{amsfonts}
\usepackage{amssymb}
\usepackage{graphicx}
\usepackage[indonesian]{babel}
\usepackage[left=2.00cm, right=2.00cm, top=2.00cm, bottom=2.00cm]{geometry}
\usepackage{float} 

\title{Tugas 10 - Pengolahan Sinyal Digital\\
		Circular Convolution}

% remove spacing around date:
\usepackage{titling}
\predate{}
\postdate{}
\date{}

\begin{document}
	\maketitle
	\date{}
	\begin{enumerate}
		\item Diketahui $ x(n) $ dan $ h(n) $ adalah dua finite length sequences, keduanya memiliki panjang $ N $.
		\begin{enumerate}
			\item Berapa panjang maksimum yang mungkin dicapai hasil linear convolution dari $ x(n) $ dengan $ h(n) $ ?
			\item Berapa panjang maksimum yang mungkin dicapai hasil N-point circula convolution dari $ x(n) $ dengan $ h(n) $ ?
		\end{enumerate}
		\item Diketahui 4-point sequences yang ditampilkan oleh Gambar \ref{fig:img01}.
		\begin{enumerate}
			\item Gambarkan linear convolution dari $ x(n) $ dengan $ x(n) $.
			\item Gambarkan 4-points circular convolution dari $ x(n) $ dengan $ x(n) $.
		\end{enumerate}
		\begin{figure}[H]
			\centering
			\includegraphics[width=0.7\linewidth]{img/img01}
			\caption{4-point sequences}
			\label{fig:img01}
		\end{figure}
	\end{enumerate}
\end{document}